\begin{figure}[ht]
\centering
\begin{tikzpicture}[x=1cm, y=1cm, font=\small]

% Plot |sin x - T_n(x)| on log scale vs x in [0, 2]
% Axes
\draw[->, thick] (-0.2, 0) -- (8.5, 0) node[right] {$\theta$ (rad)};
\draw[->, thick] (0, -0.2) -- (0, 6.2) node[above] {$|\sin\theta - T_n(\theta)|$ (log scale)};

% x-ticks at 0.25, 0.5, 1.0, 1.5, 2.0 mapped to 1, 2, 4, 6, 8 cm
\foreach \x/\lbl in {1/0.25, 2/0.5, 4/1.0, 6/1.5, 8/2.0} {
  \draw (\x, -0.1) -- (\x, 0.1) node[below, yshift=-0.15cm] {\lbl};
}
% y-ticks: log10 decades, y = 6 at 10^0, y=5 at 10^-1, etc., down to y=0 at 10^-6
\foreach \y/\lbl in {0/$10^{-6}$, 1/$10^{-5}$, 2/$10^{-4}$, 3/$10^{-3}$, 4/$10^{-2}$, 5/$10^{-1}$, 6/$10^{0}$} {
  \draw (-0.1, \y) -- (0.1, \y) node[left, xshift=-0.15cm] {\lbl};
}

% T_1 (linear): error ~ theta^3/6. log10(theta^3/6) at theta values, mapped to y = 6 + log10(...)
% theta=0.25 -> theta^3/6 = 0.0026 -> log10 = -2.58 -> y = 6 - 2.58 = 3.42; but our scale is 1 cm per decade from y=6 (10^0) down to y=0 (10^-6)
% So y = 6 + log10(value).  -2.58 -> 3.42
\draw[thick, engblue] plot coordinates {
  (0.5, 1.40)   % theta=0.125: theta^3/6 = 3.26e-4 -> log = -3.49 -> y = 2.51; but at x=0.5 mapped to theta=0.125
  (1, 2.40)     % theta=0.25: 2.60e-3 -> -2.58 -> y=3.42, but plotting steeper... use real T_1 errors below
};
% Recompute T1 error = |sin theta - theta| exact values
% theta=0.1: 1.67e-4 -> y = 6 - 3.78 = 2.22; x at 0.1 -> x_cm = 0.4
% theta=0.25: 2.60e-3 -> y = 3.42; x_cm=1
% theta=0.5: 2.06e-2 -> y = 4.31; x_cm=2
% theta=1.0: 0.1585 -> y = 5.20; x_cm=4
% theta=1.5: 0.5025 -> y = 5.70; x_cm=6
% theta=2.0: 1.091 -> y = 6.04 (clip at 6); x_cm=8

\draw[thick, engblue, mark=*, mark size=1pt] plot coordinates {
  (0.4, 2.22) (1, 3.42) (2, 4.31) (4, 5.20) (6, 5.70) (8, 6.04)
};
\node[engblue, anchor=west] at (8.1, 6.04) {$T_1$};

% T_3 (cubic): error = |sin theta - (theta - theta^3/6)|
% theta=0.1: 8.33e-8 -> y = -1.08; off-chart at bottom
% theta=0.25: 8.13e-6 -> y = 0.91; x_cm=1
% theta=0.5: 2.59e-4 -> y = 2.41; x_cm=2
% theta=1.0: 8.33e-3 -> y = 3.92; x_cm=4
% theta=1.5: 6.46e-2 -> y = 4.81; x_cm=6
% theta=2.0: 2.42e-1 -> y = 5.38; x_cm=8
\draw[thick, engred, mark=square*, mark size=1pt] plot coordinates {
  (1, 0.91) (2, 2.41) (4, 3.92) (6, 4.81) (8, 5.38)
};
\node[engred, anchor=west] at (8.1, 5.38) {$T_3$};

% T_5 (quintic): error theta=0.5: 1.55e-6 -> y=0.19; theta=1: 1.96e-4 -> y=2.29; theta=2: 0.02 -> y=4.30
\draw[thick, engblack, mark=triangle*, mark size=1pt] plot coordinates {
  (2, 0.19) (4, 2.29) (6, 3.43) (8, 4.30)
};
\node[engblack, anchor=west] at (8.1, 4.30) {$T_5$};

% Tolerance line at 1e-3
\draw[dashed, gray] (0, 3) -- (8.3, 3) node[right, gray] {$10^{-3}$};

\end{tikzpicture}
\caption[Taylor-truncation error for $\sin\theta$ at orders 1, 3, 5]{Absolute
truncation error $|\sin\theta - T_n(\theta)|$ for Taylor polynomials of
$\sin\theta$ at $0$ of degrees $1$, $3$, and $5$. Each curve has slope
$n+1$ in log-log when $\theta$ is small, matching the Lagrange remainder
bound $\theta^{n+1}/(n+1)!$. The dashed line at $10^{-3}$ shows the
budget the project's small-angle rule has to meet: $T_1$ holds it to
$\theta \approx 0.18$ rad ($10^{\circ}$); $T_3$ to $\theta \approx 1.0$
rad ($57^{\circ}$); $T_5$ comfortably to $\theta = 2$ rad. Each
additional order extends the safe range by roughly a factor of $\sqrt[3]{}$
to $\sqrt[5]{}$ of the tolerance ratio.}
\label{fig:vol02:ch04:taylor-error}
\end{figure}
